\section{선별 문제 풀이}
\label{sec:ring-selected-solutions}

\subsection*{문제 \thechapter.3: $\Z/20\Z$의 단위원과 영인수}

\begin{solution}
$\overline a$가 단위원일 필요충분조건은 $\gcd(a,20)=1$이다. 따라서
\[
  (\Z/20\Z)^\times
  =\{\overline1,\overline3,\overline7,\overline9,
      \overline{11},\overline{13},\overline{17},\overline{19}\}.
\]
역원은
\[
\begin{aligned}
&\overline1^{-1}=\overline1,
&&\overline3^{-1}=\overline7,
&&\overline7^{-1}=\overline3,
&&\overline9^{-1}=\overline9,\\
&\overline{11}^{-1}=\overline{11},
&&\overline{13}^{-1}=\overline{17},
&&\overline{17}^{-1}=\overline{13},
&&\overline{19}^{-1}=\overline{19}.
\end{aligned}
\]
영이 아닌 비단위원은 모두 영인수이다. 따라서 영이 아닌 영인수는
\[
  \overline2,\overline4,\overline5,\overline6,\overline8,
  \overline{10},\overline{12},\overline{14},\overline{15},
  \overline{16},\overline{18}
\]
이다.
\end{solution}

\subsection*{문제 \thechapter.6: 생성아이디얼}

\begin{solution}
정수환에서 두 원소가 생성하는 아이디얼은 최대공약수가 생성한다. 따라서
\[
  (18,30)=(6).
\]
두 주아이디얼의 교집합은 최소공배수가 생성하므로
\[
  (12)\cap(18)=(36).
\]
마지막으로 $X^5=X^3X^2$이므로 $(X^5)\subseteq(X^2)$이고
\[
  (X^2,X^5)=(X^2).
\]
\end{solution}

\subsection*{문제 \thechapter.8: $\operatorname{ev}_2$}

\begin{solution}
상수다항식이 모든 정수로 평가되므로 $\operatorname{ev}_2$는 전사이다. 인수정리에 의해
\[
  f(2)=0\quad\Longleftrightarrow\quad (X-2)\mid f.
\]
따라서
\[
  \ker(\operatorname{ev}_2)=(X-2),
  \qquad
  \im(\operatorname{ev}_2)=\Z.
\]
제1동형정리로
\[
  \Z[X]/(X-2)\cong\Z
\]
를 얻는다.
\end{solution}

\subsection*{문제 \thechapter.10: $\F_3[X]/(X^2+1)$}

\begin{solution}
$X^2+1$은 $\F_3$에서 근을 갖지 않으므로 기약이다. 따라서 몫환은 원소가 $9$개인 체이다. $\alpha=X+(X^2+1)$라 하면 $\alpha^2=-1=2$이고 원소는
\[
  0,1,2,\alpha,2\alpha,1+\alpha,1+2\alpha,
  2+\alpha,2+2\alpha
\]
이다. 영이 아닌 원소의 역원은
\[
\begin{aligned}
&1^{-1}=1,\qquad 2^{-1}=2,\qquad
\alpha^{-1}=2\alpha,\qquad(2\alpha)^{-1}=\alpha,\\
&(1+\alpha)^{-1}=2+\alpha,\qquad
(2+\alpha)^{-1}=1+\alpha,\\
&(1+2\alpha)^{-1}=2+2\alpha,\qquad
(2+2\alpha)^{-1}=1+2\alpha.
\end{aligned}
\]
일반적으로 $a+b\alpha\ne0$의 역원은
\[
  (a+b\alpha)^{-1}=\frac{a-b\alpha}{a^2+b^2}
\]
로 계산할 수 있다.
\end{solution}

\subsection*{문제 \thechapter.12: 멱영원소들의 아이디얼}

\begin{solution}
$0$은 멱영원소이다. $a^m=0$이면 $(-a)^m=(-1)^ma^m=0$이므로 덧셈 역원에 닫혀 있다.

$a^m=0$, $b^n=0$이라 하자. 가환성을 이용해 $(a+b)^{m+n-1}$을 이항전개하면 각 항은
\[
  \binom{m+n-1}{k}a^kb^{m+n-1-k}
\]
꼴이다. $k\ge m$이면 $a^k=0$이고, $k<m$이면 $m+n-1-k\ge n$이므로 $b^{m+n-1-k}=0$이다. 따라서 모든 항이 $0$이고 $a+b$도 멱영원소이다.

마지막으로 $r\in R$이면
\[
  (ra)^m=r^ma^m=0
\]
이므로 $ra$도 멱영원소이다. 따라서 멱영원소 전체는 아이디얼이다.
\end{solution}

\subsection*{문제 \thechapter.14: 아이디얼의 역상}

\begin{solution}
$a,b\in\varphi^{-1}(J)$이면 $\varphi(a),\varphi(b)\in J$이고
\[
  \varphi(a-b)=\varphi(a)-\varphi(b)\in J.
\]
또한 $r\in R$에 대해
\[
  \varphi(ra)=\varphi(r)\varphi(a)\in J
\]
이므로 $ra\in\varphi^{-1}(J)$. 따라서 역상은 아이디얼이다.

$J$가 소이고 $ab\in\varphi^{-1}(J)$이면
\[
  \varphi(a)\varphi(b)=\varphi(ab)\in J.
\]
소아이디얼의 정의에 의해 $\varphi(a)\in J$ 또는 $\varphi(b)\in J$이므로 $a$ 또는 $b$가 역상에 속한다. 또한 $1_S\notin J$이고 $\varphi(1_R)=1_S$이므로 $1_R\notin\varphi^{-1}(J)$이다. 따라서 역상은 진아이디얼이며 소아이디얼이다.

마지막으로 $\varphi$가 전사이고 $J$가 극대라 하자. 합성
\[
  R\xrightarrow{\varphi}S\twoheadrightarrow S/J
\]
은 전사이고 그 핵은 $\varphi^{-1}(J)$이다. 따라서 제1동형정리에 의해
\[
  R/\varphi^{-1}(J)\cong S/J.
\]
오른쪽이 체이므로 왼쪽도 체이고, 따라서 $\varphi^{-1}(J)$는 극대아이디얼이다.
\end{solution}

\subsection*{문제 \thechapter.16: 제2동형정리}

\begin{solution}
먼저 $1\in S\subseteq S+J$이고, 차와 곱에 대한 닫힘성을 직접 확인하면 $S+J$가 $R$의 부분환임을 알 수 있다. 또한 $S\cap J$는 $S$의 덧셈부분군이며 $s\in S$, $x\in S\cap J$이면 $sx\in S\cap J$이므로 $S$의 아이디얼이다.

이제 사상
\[
  \psi:S\to(S+J)/J,
  \qquad s\mapsto s+J
\]
를 생각하자. 임의의 $(s+j)+J$에서 $j\in J$이므로
\[
  (s+j)+J=s+J
\]
이고, 따라서 $\psi$는 전사이다. 또한
\[
  \ker\psi=S\cap J.
\]
제1동형정리에 의해
\[
  S/(S\cap J)\cong(S+J)/J
\]
이다.
\end{solution}

\subsection*{문제 \thechapter.18: 아이디얼의 거듭제곱도 서로소}

\begin{solution}
$I+J=R$이므로 $x\in I$, $y\in J$ 중 $x+y=1$인 것이 있다. 이항정리를 적용하면
\[
  1=(x+y)^{m+n-1}
   =\sum_{k=0}^{m+n-1}
     \binom{m+n-1}{k}x^ky^{m+n-1-k}.
\]
각 항에서 $k\ge m$이면 그 항은 $I^m$에 속한다. $k<m$이면
\[
  m+n-1-k\ge n
\]
이므로 그 항은 $J^n$에 속한다. 따라서 $1\in I^m+J^n$이고
\[
  I^m+J^n=R.
\]
\end{solution}

\subsection*{문제 \thechapter.20: 체의 세 가지 특징}

\begin{solution}
$R\ne0$라 하자. $R$이 체이면 준동형 $\varphi:R\to S$의 핵은 $R$의 아이디얼이다. $S\ne0$이므로 $\varphi(1)=1\ne0$이고 핵은 $R$ 전체가 아니다. 체의 유일한 진아이디얼은 $(0)$이므로 $\varphi$는 단사이다.

(ii)를 가정하고 몫사상 $R\to R/I$를 생각하자. $I$가 진아이디얼이면 $R/I\ne0$이므로 몫사상은 단사이다. 그 핵이 $I$이므로 $I=(0)$이다. 따라서 아이디얼은 서로 다른 두 아이디얼 $(0)$과 $R$뿐이다.

(iii)을 가정하자. $0\ne a\in R$이면 $(a)$는 영아이디얼이 아니므로 $(a)=R$이다. 따라서 $1\in(a)$이고 어떤 $b\in R$에 대해 $ab=1$이다. 모든 영이 아닌 원소가 단위원이므로 $R$은 체이다.
\end{solution}

\subsection*{문제 \thechapter.21: 세 합동식}

\begin{solution}
먼저 $x\equiv3\pmod{10}$이라 두어 $x=3+10k$로 쓴다. 첫째 합동식은
\[
  3+k\equiv4\pmod9
\]
이므로 $k\equiv1\pmod9$이다. 따라서 $k=1+9t$이고
\[
  x=13+90t.
\]
셋째 합동식에 대입하면
\[
  13+90t\equiv6+6t\equiv5\pmod7.
\]
따라서 $6t\equiv6\pmod7$, 즉 $t\equiv1\pmod7$이다. 결론적으로
\[
  x\equiv103\pmod{630}.
\]
\end{solution}

\subsection*{문제 \thechapter.23: $K[X]/(X^2-1)$}

\begin{solution}
사상
\[
  \Phi:K[X]\to K\times K,
  \qquad f\mapsto(f(1),f(-1))
\]
를 생각하자. 표수가 $2$가 아니므로 임의의 $(a,b)\in K\times K$에 대해
\[
  f(X)=\frac{a+b}{2}+\frac{a-b}{2}X
\]
라 두면 $f(1)=a$, $f(-1)=b$이다. 따라서 $\Phi$는 전사이다. 핵은 $1$과 $-1$을 모두 근으로 갖는 다항식들이므로
\[
  \ker\Phi=(X-1)(X+1)=(X^2-1).
\]
제1동형정리로 원하는 동형을 얻는다. 특히
\[
  X+(X^2-1)\longmapsto(1,-1).
\]
\end{solution}

\subsection*{문제 \thechapter.25: $\Q[X]/(X^3-2)$}

\begin{solution}
$\alpha^3=2$를 이용하면
\[
  (1+\alpha+\alpha^2)(1-\alpha)
  =1-\alpha^3=-1.
\]
또한
\[
  (1+\alpha+\alpha^2)(\alpha-1)
  =\alpha^3-1=1
\]
이므로
\[
  (\alpha^2+\alpha+1)^{-1}=\alpha-1.
\]
다항식 $X^3-2$는 유리근을 갖지 않는다. 삼차다항식은 가약이면 일차인수를 가지므로 $X^3-2$는 $\Q[X]$에서 기약이다. 따라서 $(X^3-2)$는 극대아이디얼이고 $A$는 체이다.
\end{solution}

\subsection*{문제 \thechapter.27: PID의 최대공약수}

\begin{solution}
$(a,b)=(d)$이면 $a,b\in(d)$이므로 $d\mid a$, $d\mid b$이다. 한편 $d\in(a,b)$이므로 어떤 $x,y\in R$에 대해
\[
  d=xa+yb
\]
이다. 따라서 $a,b$의 모든 공약수는 $d$를 나눈다. 즉 $d$는 최대공약수이다.

일반 UFD에서는 최대공약수가 존재하지만 $(a,b)$를 생성하지 않을 수 있다. 예를 들어 $K[X,Y]$는 UFD이고 $\gcd(X,Y)=1$이지만
\[
  (X,Y)\ne(1)=K[X,Y].
\]
실제로 $(X,Y)$의 모든 다항식은 상수항이 $0$이므로 $1\notin(X,Y)$이다.
\end{solution}

\subsection*{문제 \thechapter.28: 가우스 정수의 몫환}

\begin{solution}
직접 곱하면
\[
  (2+i)(2-i)=4-i^2=5
\]
이다. 가우스 노름 $N(a+bi)=a^2+b^2$는 곱셈적이다. 단위원 $u$는
\[
  1=N(1)=N(u)N(u^{-1})
\]
을 만족하므로 $N(u)=1$이어야 한다. 그런데 $N(2+i)=5$이므로 $2+i$는 단위원이 아니다.

이제 환 준동형
\[
  \varphi:\Z[i]\to\F_5,
  \qquad a+bi\mapsto\overline{a+3b}
\]
를 정의하자. $3^2=9\equiv-1\pmod5$이므로 곱셈을 보존하며, 정수들의 상만으로도 $\F_5$ 전체에 도달하므로 전사이다. 또한
\[
  \varphi(2+i)=\overline{2+3}=0
\]
이어서 $(2+i)\subseteq\ker\varphi$이다.

반대로 $a+bi\in\ker\varphi$이면 $a+3b\equiv0\pmod5$, 즉 $a\equiv2b\pmod5$이다. 따라서
\[
  5\mid(2a+b),\qquad 5\mid(2b-a),
\]
이고
\[
  a+bi=(2+i)\left(\frac{2a+b}{5}+\frac{2b-a}{5}i\right).
\]
괄호 안은 다시 가우스 정수이므로 $a+bi\in(2+i)$이다. 따라서 $\ker\varphi=(2+i)$이고 제1동형정리에 의해
\[
  \Z[i]/(2+i)\cong\F_5.
\]
이 동형 아래에서 $i$는 $\overline3=\overline{-2}$에 대응하며
\[
  \overline3^{\,2}=\overline9=\overline4=\overline{-1}
\]
이다.
\end{solution}

\subsection*{문제 \thechapter.29: 소아이디얼과 극대아이디얼}

\begin{solution}
(a) 모든 정역에서 $(0)$은 소아이디얼이다. 만약 모든 소아이디얼이 극대라면 $(0)$도 극대이고
\[
  R/(0)\cong R
\]
가 체가 된다.

(b) $R$이 PID이고 $(p)\ne(0)$가 소아이디얼이라고 하자. 그러면 $p$는 소원소이고 따라서 기약원소이다. PID에서는 기약원소가 생성하는 아이디얼이 극대이므로 $(p)$는 극대아이디얼이다.

(c) $\Z$는 PID이므로 모든 영이 아닌 소아이디얼 $(p)$가 극대이다. 그러나 $(0)$도 소아이디얼이면서 극대는 아니다. 실제로
\[
  (0)\subsetneq(2)\subsetneq\Z.
\]
따라서 $\Z$는 (b)의 결론은 만족하지만 ``모든 소아이디얼이 극대''라는 (a)의 가정은 만족하지 않는다.
\end{solution}
